\scanpage{72}
C'est \unsure{aussi} équivalent au pb suivant : si $X$ est comme dessus, et $X_0-X$ une imm. \uline{nilp. loc.}, $\unsure{Y|X}$ de prés. finie sur $S$, a-t-on
\[
H^*(X_{\unsure{\mathrm{strat}}},F)\xrightarrow{\sim}H^*(X_{0\,\mathrm{str.}},F)
\]
[De plus, on peut supposer $S,X,X_0$ affines].

\uline{Remarque} Si $S=\operatorname{Spec}k$, $k$ corps valué complet non archimédien, il y a lieu de \deleted{faire} se \deleted{\unsure{demander}}\inserted{poser la question de la validité de la} conjecture \unsure{plus} ci-dessus, pour les \deleted{espaces \uline{rigide-analytiques} de Tate (au sens \unsure{strict}, [\unsure{et rigides}, --- mais en \unsure{considérant} les} espaces analytiques \deleted{\unsure{ordinaires}} ``flasques'' sur $k$ ].
\deleted{Cela expliquerait} Mais cela n'aurait pas d'implications globales si remarquables ....
\mspage{71}
